\documentclass{article}

\usepackage{layout}
\setlength{\textwidth}{15.0 cm}
\setlength{\textheight}{25.0 cm}


\usepackage[english,brazil]{babel}
\usepackage{pagina}	% pagina-padrao
\usepackage{indentfirst}		% for indent
\usepackage[utf8]{inputenc}
\usepackage{graphics,epsfig}
\usepackage{graphics}
\graphicspath{{./figuras/}}
\usepackage{pstricks,pst-node,pst-tree}
\usepackage{alltt}
%\usepackage{makeidx}
%\makeindex
\usepackage[figuresright]{rotating} % for saydways tables and figures
\usepackage{enumerate}			% for configuration of enumerate environment
\usepackage{amsmath}
\usepackage{amssymb}
\usepackage{portland,multirow}
\usepackage{float}
\usepackage{graphicx}


\setcounter{secnumdepth}{3}	% numeracao ate subsubsecao
\setcounter{tocdepth}{2}	% indice ate subsubsecao

\usepackage{longtable}


\begin{document}

\begin{center}
    
\textbf{UNIVERSIDADE FEDERAL DO RIO DE JANEIRO}
\vspace{-0.2cm}

\textbf{ESCOLA POLITÉCNICA}
\vspace{-0.2cm}

\textbf{DEPARTAMENTO DE ENGENHARIA ELETRÔNICA E DE COMPUTAÇÃO}
\vspace{0.8cm}

\underline{\textbf{PROPOSTA DE PROJETO DE GRADUAÇÃO}}

Aluno: Rafael Magalhães de Matos Gabriel
\vspace{-0.2cm}

rafaeldzn2003@poli.ufrj.br

Orientador: Fernando Antônio Pinto Barúqui

\vspace{0.2cm}

Curso: (x) Eletrônica e Computação;  ( ) Computação e Informação.

\end{center}

\textbf{1. TÍTULO}

    Amplificador de Potência Estéreo Classe D de 20 Watts RMS para Áudio

\vspace{0.4cm}

\textbf{2. ÊNFASE}

    Eletrônica

\vspace{0.4cm}

\textbf{3. TEMA}

O presente trabalho apresenta o desenvolvimento de um amplificador de potência estéreo para áudio, na configuração classe D,
que visa alta eficiência energética e baixo THD, para atender as normas de alta fidelidade.

% Duvidas: eu consigo fazer um classe D com baixo THD? Colocar como objetivo a alta eficiência energética faz sentido?

\end{document}